\begingroup
\setlength{\tabcolsep}{3.5pt}
\begin{table}[h]
    \centering
    \begin{tabular}{ll cccc cccc cc}
    \toprule
     & & \multicolumn{8}{c}{\textsc{Prompt Tuning Analysis}} \\
     \cmidrule(lr){3-10} 
     
    & \makecell[c]{\scriptsize{Prompt tuning}\vspace{-1.2mm}\\\scriptsize{train.~examples}} & \datasetacc{BoolQ} & \datasetacc{CB} & \datasetacc{CoPA} & \datasetfone{MultiRC} & \datasetacc{ReCoRD} & \datasetacc{RTE} & \datasetacc{WiC} & \datasetacc{WSC}  \\
     \midrule
     \baselm{} & \multirow{2}{*}{32} & 55.5 & 55.4 & 87.0 & 65.4 & 78.0 & 52.4 & 51.6 & 65.4 \\
     \flan{} & & 77.5 & 87.5 & 91.0 & 76.8 & 80.8 & 83.0 & 57.8 & 70.2 \\
     \midrule
     \baselm{} & \multirow{2}{*}{\makecell[l]{\footnotesize{full}\vspace{-0.5mm}\\\footnotesize{dataset}}} & 82.8 & 87.5 & 90.0 & 78.6 & 84.8 & 82.0 & 54.9 & 72.7 \\
     \flan{} & & 86.3 & 98.2 & 94.0 & 83.4 & 85.1 & 91.7 & 74.0 & 86.5 \\
    \bottomrule
    \end{tabular}
    \caption{
    \flan{} (instruction tuning) responds better to continuous inputs attained via prompt tuning than \baselm{} (no instruction tuning).
    When prompt tuning on a given dataset, no tasks from the same cluster as that dataset were seen during instruction tuning.
    }
    \label{tab:prompt_tuning}
\end{table}
\endgroup